% ===== Document Preamble =====
\documentclass[11pt]{article}
\usepackage[margin=0.5in]{geometry}
\usepackage{amsmath,amssymb}
\usepackage{xcolor}
\usepackage{enumitem}
\usepackage{needspace}
\usepackage{comment}

% ===== Colors =====
\definecolor{AnsRed}{RGB}{170,0,0}
\definecolor{WrongBlue}{RGB}{0,0,160}
\definecolor{SkipGray}{RGB}{120,120,120}

% ===== Choice list styling (tight, consistent) =====
\setlist[enumerate]{%
  leftmargin=1.5em,
  itemsep=0.25em,
  topsep=0.25em,
  label=\textbf{\Alph*.}
}

% ===== Style Wrappers =====
\newcommand{\QuestionStyle}[1]{{\large\bfseries\boldmath #1}}
\newcommand{\AnswerStyle}[1]{{\large\bfseries\boldmath #1}}
\newcommand{\ReasonStyle}[1]{{\normalsize #1}}

% ===== Show / Hide Reasoning Toggle =====
\newif\ifshowreason
\showreasontrue      % Instructor version
% \showreasonfalse   % Student version

% ===== Choice Item Macros (Answer + Reasoning) =====
% Usage:
%   \Correct{Answer}{Reasoning}
%   \Wrong{Answer}{Reasoning}

\newcommand{\Correct}[2]{%
  \item \AnswerStyle{\textcolor{AnsRed}{#1}}%
  \ifshowreason\\[-0.2em]
    \ReasonStyle{\textcolor{AnsRed}{#2}}%
  \fi
}

\newcommand{\Wrong}[2]{%
  \item \AnswerStyle{\textcolor{WrongBlue}{#1}}%
  \ifshowreason\\[-0.2em]
    \ReasonStyle{\textcolor{WrongBlue}{#2}}%
  \fi
}

% ===== Multiple-Choice Question Block Macro =====
% Usage: \QBlock{Q1}{stem...}{<items>}{answer key}

\newcommand{\QBlock}[4]{%
  \Needspace{8\baselineskip}%
  \noindent
  \begin{minipage}{\linewidth}
    \QuestionStyle{#1.~#2}

    \vspace{0.25em}
    \begin{enumerate}
      #3
    \end{enumerate}

    \vspace{0.1em}
    \noindent\textcolor{AnsRed}{\textit{\textbf{Answer:}~#4}}

    \vspace{0.25em}
    \hrule height 0.4pt
  \end{minipage}
  \vspace{0.75em}
}

% ===== Open-Ended Question Block Macro =====
% Usage: \OEBlock{Q1}{stem...}{answer text}

\newcommand{\OEBlock}[3]{%
  \Needspace{5\baselineskip}%
  \noindent
  \begin{minipage}{\linewidth}
    \QuestionStyle{#1.~#2}

    \vspace{0.25em}
    \noindent\textcolor{AnsRed}{\textit{\textbf{Answer:}~#3}}

    \vspace{0.25em}
    \hrule height 0.4pt
  \end{minipage}
  \vspace{0.75em}
}

% ===== Multi-Part Question Block Macro =====
% Usage: \MPBlock{Q1}{stem...}{<parts>}

\newcommand{\MPBlock}[3]{%
  \Needspace{8\baselineskip}%
  \noindent
  \begin{minipage}{\linewidth}
    \QuestionStyle{#1.~#2}

    \vspace{0.25em}
    #3

    \vspace{0.25em}
    \hrule height 0.4pt
  \end{minipage}
  \vspace{0.75em}
}

% ===== MC Sub-Part Macro =====
% Usage: \MCPart{a}{stem...}{<items>}{answer key}

\newcommand{\MCPart}[4]{%
  \par\vspace{0.35em}
  \noindent\textbf{(#1)}~#2
  \begin{enumerate}
    #3
  \end{enumerate}
  \vspace{0.1em}
  \noindent\textit{\textbf{Answer:}~#4}
  \par
}

% ===== Open-Ended Sub-Part Macro =====
% Usage: \OEPart{a}{stem...}{answer text}

\newcommand{\OEPart}[3]{%
  \par\vspace{0.35em}
  \noindent\textbf{(#1)}~#2
  \par\vspace{0.15em}
  \noindent\textit{\textbf{Answer:}~#3}
  \par
}

% ===== Skipped Question Note Macro =====
% Usage: \SkipBlock{Q1}{reason}

\newcommand{\SkipBlock}[2]{%
  \noindent\textcolor{SkipGray}{\textit{[#1 omitted: #2]}}%
  \vspace{0.5em}\par
}

% ===== Document Start =====

\begin{document}

\begin{center}
{\huge \textbf{Quiz \#2}}\\[0.25em]
{\Large physics\_abstract}\\[0.25em]
{\large Questions generated below. Multiple-choice answers in RED; incorrect answers in BLUE.}
\end{center}

\vspace{0.75em}

\QBlock{Q1}{A particle is located relative to the origin at the position vector $\vec r=-7\hat{i}+11\hat{j}$, where distances are measured in arbounits. Which of the following statements is correct?}{
\Wrong{The distance from the particle to the origin is 18 arbounits.}{This comes from adding the components, $7+11=18$, instead of finding the magnitude of the vector. Distance from the origin must be computed with the Pythagorean relation, not by direct addition.}
\Wrong{The distance from the particle to the origin is 11 arbounits.}{This mistake treats the $y$-component as if it were the full distance. It ignores the nonzero $x$-component of the position vector.}
\Wrong{The distance from the particle to the origin is 7 arbounits.}{This mistake treats the $x$-component magnitude as the total distance. It neglects the contribution from the $y$-component.}
\Correct{None of the above is correct.}{The distance from the origin is the magnitude of the position vector, $|\vec r|=\sqrt{(-7)^2+11^2}=\sqrt{170}\approx 13.0$ arbounits. Since $13.0$ is not equal to $18$, $11$, or $7$, the only correct choice is that none of the listed numerical values is correct.}
}{D}

\QBlock{Q2}{If the displacement of an object is given by $\vec s(t)=\big((2.1\,\text{m/s}^3)t^3-(1.7\,\text{m/s}^2)t^2\big)\hat{i}$, which of the following statements is correct?}{
\Wrong{This is a constant-acceleration motion.}{This assumes the acceleration does not change with time, but differentiating a cubic displacement produces a time-dependent acceleration. The presence of the $t^3$ term guarantees that $\vec a(t)$ is not constant.}
\Correct{The magnitude of the acceleration at $t=0$ is $3.4\,\text{m/s}^2$.}{Differentiate the displacement twice: $\vec v(t)=\big(6.3t^2-3.4t\big)\hat{i}$ and $\vec a(t)=\big(12.6t-3.4\big)\hat{i}$. At $t=0$, the acceleration is $-3.4\hat{i}\,\text{m/s}^2$, so its magnitude is $3.4\,\text{m/s}^2$.}
\Wrong{The magnitude of the acceleration is $1.7\,\text{m/s}^2$.}{This confuses the coefficient of the $t^2$ term in displacement with the acceleration itself. After two derivatives, that coefficient contributes $-3.4\,\text{m/s}^2$, not $1.7\,\text{m/s}^2$.}
\Wrong{The magnitude of the acceleration at $t=0$ is $0\,\text{m/s}^2$.}{This error comes from focusing only on the cubic term, whose contribution to acceleration vanishes at $t=0$. It overlooks the quadratic term, whose second derivative gives a nonzero constant contribution.}
}{B}

\QBlock{Q3}{An object takes $4.0\,\text{s}$ to move from position $\vec r_A=1.0\hat{i}-2.0\hat{j}$ to position $\vec r_B=7.0\hat{i}+6.0\hat{j}$, where position is measured in meters. Which of the following statements is correct?}{
\Wrong{The path length traveled by the object is $10\,\text{m}$.}{This uses the displacement magnitude, $\sqrt{(7-1)^2+(6-(-2))^2}=10\,\text{m}$, as though it were the actual distance traveled along the path. The path length depends on the trajectory and cannot be determined from just the initial and final positions.}
\Wrong{The speed at $t=0$ must be $2.5\,\text{m/s}$.}{This comes from dividing the displacement magnitude by the total time to get $10/4=2.5\,\text{m/s}$, then incorrectly interpreting that as the initial speed. That calculation gives neither the initial speed nor necessarily even the average speed along the actual path.}
\Correct{None of the above is correct.}{From the given data, we can determine only the displacement, $\Delta\vec r=(6\hat{i}+8\hat{j})\,\text{m}$, whose magnitude is $|\Delta\vec r|=10\,\text{m}$. However, the actual path length, the speed at $t=0$, and the acceleration at $t=0$ are not uniquely determined by just two positions and one elapsed time, so the other statements are not necessarily true.}
\Wrong{The acceleration at $t=0$ must be $0.63\,\text{m/s}^2$.}{This assigns a specific acceleration value from insufficient information. Infinitely many possible motions can connect the same start and end positions in $4.0\,\text{s}$, each with different accelerations.}
}{C}

\end{document}
